\section{Engineering und Testing (Kap. 05)}
\label{sec:engineering-testing}

\subsection{Framework- und Implementierungsentscheidungen}

Die MAS-Laufzeit ist bewusst ohne zusätzliches Agentenframework implementiert.
\texttt{MASAgent} definiert den gemeinsamen Agentenvertrag,
\texttt{MASMessageBus} eine FIFO-Mailbox mit sanitisiertem Kommunikations-Trace
und \texttt{AgentBlackboard} den gemeinsamen Laufzustand. Der Orchestrator
kennt Agentenregister, Nachrichtenwarteschlange und Abbruchgrenze, aber nicht
die fachliche Entscheidung: Das Routing nach Kritik erzeugt der Reviewer.
Pydantic validiert Pläne, Reviews, Nachrichten, Toolstatus, lokale Zustände und
API-Antworten. Sämtliche Laufobjekte werden je Request erzeugt. Dadurch teilt
der globale FastAPI-Orchestrator keinen veränderlichen Agentenzustand zwischen
parallelen Nutzern.

Diese Eigenimplementierung hält den kleinen Zustandsgraphen und die
Capability-Policy direkt prüfbar. Ein graphorientiertes Framework wäre bei
persistenten Checkpoints, vielen Verzweigungen oder Human-in-the-Loop-Pausen
neu abzuwägen. Gegenwärtig würde es zusätzliche Abhängigkeit und Migration
erzeugen, ohne eine benötigte Funktion zu ersetzen. Der Preis des eigenen
Runtimes bleibt Wartungsaufwand: Zustandsmigration, Wiederaufnahme und
verteilte Ausführung müssten selbst entwickelt werden. Das System ist
asynchron programmiert, dispatcht die Agenten aber absichtlich seriell.

Der OpenAI-kompatible Adapter kapselt JSON-Modus und Function Calling. Die
Tool-Registry trennt Schema und Implementierung, während
\texttt{SAFE\_TOOL\_POLICY} Rechte aktionsabhängig reduziert. Toolresultate
werden als Erfolg, Fehler, Blockade oder Argumentfehler typisiert. Der Reviewer
kann auf dieser Basis Textrevision oder Neuplanung auslösen. Ein Limit von eins
bis vier, standardmäßig zwei Koordinationsrunden verhindert Endlosschleifen.
Weiterhin fehlen persistente Checkpoints, Retries mit Backoff und ein
transaktionales Zurückrollen bereits ausgeführter Tools.

\subsection{Testfälle und Beobachtungen}

Alle 64 Tests bestanden. Die einzige Warnung betrifft die künftige
Umstellung des Starlette-Testclients. Zusätzlich bestanden zehn Angular-Tests
und der Angular-Produktionsbuild. Beide Docker-Compose-Profile wurden
erfolgreich aufgelöst. Die drei in \Cref{tab:tests} ausgewählten Tests
adressieren gerade die Merkmale, die eine Rollen-Pipeline nicht belegen könnte.

Ein ergänzender Live-End-to-End-Lauf mit dem konfigurierten OpenAI-Modell
verarbeitete ein zweiseitiges Test-PDF zu vier Aufgaben, einem Quiz mit sechs
Fragen und einem Merkblatt. Der Hinweis-Turn beteiligte Planner, Tutor und
Reviewer über sechs Nachrichten und zwei Runden. Im Bewertungsturn wurde
\texttt{evaluate\_answer} mit Status \texttt{succeeded} protokolliert. Eine
dokumentgestützte Multiple-Choice-Übung wurde ebenfalls erzeugt. Der Lauf ist
ein Integrationsnachweis, aber keine Nutzerstudie.

\begin{table}[htbp]
  \centering
  \footnotesize
  \begin{tabularx}{\linewidth}{@{}p{0.25\linewidth}XX@{}}
    \toprule
    Testfall & Prüfkriterium & Beobachtung \\
    \midrule
    Nachrichten und lokaler Zustand & Erwartet werden die typisierte Folge
    \texttt{plan\_request}, \texttt{execute\_plan},
    \texttt{review\_request}, \texttt{final\_response} sowie je eine lokale
    Entscheidung der drei Agenten. & Test bestanden. Die Empfängerfolge war
    Planner, Tutor, Reviewer, Coordinator. \\
    Reviewer-Revision & Ein Entwurf verrät die Lösung. Der Reviewer muss Kritik
    an den Tutor senden. Die toolfreie Revision wird erneut geprüft. & Test
    bestanden. Es gab zwei Runden, zwei Tutor- und zwei Reviewer-Entscheidungen. Die
    textbezogene Assertion prüft den Beginn der ausgelieferten Frage. \\
    Neuplanung nach Toolfehler & Eine fehlgeschlagene
    \texttt{evaluate\_answer}-Beobachtung muss den Reviewer zu
    \texttt{replan\_request} und den Planner zu einem neuen
    \texttt{clarify}-Plan veranlassen. & Test bestanden. Der zweite Plan und
    die Fehlerbeobachtung waren im Laufzustand nachweisbar. \\
    \bottomrule
  \end{tabularx}
  \caption{Ausgewählte, tatsächlich ausgeführte MAS-Tests.}
  \label{tab:tests}
\end{table}

Die markerbasierten semantischen Assertions des kontrollierten Experiments
nutzen vorgegebene Modellantworten, aber den produktiven Tool- und
MAS-Kontrollfluss. Die
ReAct-Variante erreichte bei drei konstruierten Läufen Score und Tool-Abdeckung
von 1,00. Die handgefertigte Baseline erreichte 0,27 beziehungsweise 0,00. Ein zusätzlicher
MAS-Lauf erzeugte reproduzierbar \texttt{plan\_request} $\rightarrow$
\texttt{execute\_plan} $\rightarrow$ \texttt{review\_request} $\rightarrow$
\texttt{revision\_request} $\rightarrow$ \texttt{review\_request}
$\rightarrow$ \texttt{final\_response} und endete nach zwei Runden mit einer
sokratischen Frage. Dies belegt
Nachrichtenrouting und Rückkopplung, nicht fachliche Modellqualität.

Aussagen zu Tutorqualität und Lernfortschritt benötigen deshalb Live-LLM- und
Nutzerstudien. Die Angular-Suite prüft UI-Komponenten. Jede neue Tutorantwort
zeigt außerdem eine aufklappbare Zusammenfassung aus beteiligten
Agenten, Rundenzahl, Nachrichtentypen und Toolstatus. Vollständige Payloads
bleiben auf die interne Laufzeit beschränkt.
